\scanpage{44}
\uline{Foncteurs adjoints}

\uline{Définition 1} soient $X,Y\in(\mathrm{Cat})$, $f:X\longrightarrow Y$ et $g:Y\longrightarrow X$, on dit que $(f,g)$ est un couple de foncteurs adjoints, ou que $g$ est adjoint de $f$, ou que $f$ est coadjoint de $g$, s'il existe un

\deleted{$\unsure{\forall}$}
\figure{\begin{tikzpicture}[x=1cm,y=1cm]
\node[anchor=west] at (.3,2.5) {$\mathcal H^1(X,\underline{\operatorname{Hom}}(\gamma,G))\simeq$};
\node[anchor=west] at (.3,1.7) {$\mathcal H^1(X,\underline{\operatorname{Ext}}^1(\gamma,G))$};
\node[anchor=west] at (1,.45) {$\operatorname{Ext}^1(\gamma,$};
\draw (0,0)--(3.5,.12)--(6,4.1)--(10,4.05);
\draw (0,.27)--(3.5,.12);
\end{tikzpicture}}
\diagram{\begin{tikzcd}[ampersand replacement=\&, row sep=large]
X\arrow[r]\&\underline{\operatorname{Hom}}_{S-gr}(\gamma,G)\\
X\times\gamma\arrow[r]\&G\\
\deleted{\unsure{\gamma}}\ \gamma\arrow[r]\arrow[dr,no head]\&\underline{\operatorname{Hom}}_S(X,G)\\
S\&\overset{\unsure{\Pi_2}}{C}
\end{tikzcd}}
\[\operatorname{Hom}(\gamma,\]
\[\Pi^0_{X/S}\gamma_X\]
\[\text{''}\]
\[\underline{\operatorname{Hom}}_S(X,\gamma)\unsure{=|-}\gamma\]
\diagram{\begin{tikzcd}[ampersand replacement=\&, column sep=small]
\unsure{e|0}\arrow[r]\&\gamma\arrow[r]\&\underline{\operatorname{Hom}}_S(X,\gamma)\arrow[r]\&\operatorname{Hom}(X,a;\gamma)\arrow[r]\&0
\end{tikzcd}}
\figure{\begin{tikzpicture}\draw (0,0)--(.35,.12);\end{tikzpicture}}
\margin{43}
